\section{The collision of a pair and its complementary divisor}
\label{sec:isolated-pair}

This section proves a local collision formula which keeps every complementary zero and the full real drift. Throughout, $D$ denotes differentiation in the spatial variable, and negative heat time is written $t=-a$, so $f_{-a}=e^{aD^2}f$.

\begin{lemma}[A centered canonical product]\label{lem:centered-product}
Let $G$ be real entire of order at most one, with real zeros and $G(b)\ne0$ for $b\in\mathbb R$. Put
\[
 L=G'(b)/G(b),\quad S=\sum_{G(\xi)=0}(b-\xi)^{-2},\quad
 R(h)=G(b+h)/(G(b)e^{Lh}),
\]
where the zeros are repeated according to multiplicity. Then
\[
 R(h)=\prod_{G(\xi)=0}
       \left(1+\frac h{b-\xi}\right)e^{-h/(b-\xi)},
 \qquad |R(h)|\le e^{S|h|^2/2}.
\]
Moreover, with $T_3=\sum_{G(\xi)=0}(b-\xi)^{-3}$,
\[
 R(h)=1-\tfrac12Sh^2+\tfrac13T_3h^3+E_4(h),\qquad
 |T_3|\le S^{3/2},\quad
 |E_4(h)|\le C S^2|h|^4e^{2S|h|^2},
\]
for an absolute constant $C$. Also
$|R(h)-1|\le8S|h|^2e^{2S|h|^2}$.
\end{lemma}
\begin{proof}
The genus-one Hadamard product has at most a linear exponential factor. Centering it at $b$ and matching the value and logarithmic derivative gives the product written above. For $u\in\mathbb C$, the inequality
$\log(1+v)\le v$ with $v=2\Re u+|u|^2\ge-1$ gives
$|(1+u)e^{-u}|\le e^{|u|^2/2}$; a zero factor satisfies the same inequality directly. Multiplying and taking the locally uniform limit proves the growth bound.

The expansion of $\log R$ at zero gives the quadratic and cubic coefficients. The inequality
$\sum |b-\xi|^{-3}\le(\sum |b-\xi|^{-2})^{3/2}$ proves the stated bound for $T_3$.
For $S>0$ put $v=h\sqrt S$. On $|v|=1$ the normalized function is bounded by $e^{1/2}$. Cauchy's coefficient bound and the geometric series bound the remainder after degree three by $2e^{1/2}|v|^4$ on $|v|\le1/2$. Outside that disk, the global growth estimate and the three explicit coefficients are bounded by a fixed multiple of $|v|^4e^{2|v|^2}$. The same argument after degree one gives the claimed bound with constant $8$. For $S=0$ the product is empty and both remainders vanish.
\end{proof}

\begin{theorem}[An isolated pair]\label{thm:isolated-pair}
Let $f$ be real entire of order at most one with real zeros, and let $b-\delta/2,b+\delta/2$ be adjacent simple zeros, where $\delta>0$. Define
\[
 G(z)=\frac{f(z)}{(z-b)^2-\delta^2/4},\quad
 L=G'(b)/G(b),\quad S=\sum_{G(\xi)=0}(b-\xi)^{-2},\quad
 \kappa=\delta^2S.
\]
There is an absolute $\kappa_0>0$ such that every pair with $\kappa<\kappa_0$ has exactly one collision of its two continued zeros during $0\le a\le\delta^2/4$. This collision is double, and
\[
 a_c=\delta^2\left(\frac18+\frac\kappa{32}+O(\kappa^2)\right),
 \qquad z_c=b-2a_cL+O(\delta\kappa^{3/2}).
\]
The constants in these estimates are absolute. The two zeros are real before $a_c$ and form a nonreal conjugate pair afterwards on the stated interval.
\end{theorem}
\begin{proof}
Set $a=\delta^2\tau$ and $z=b-2aL+\delta w$. The exact operator identity
\[
 e^{aD^2}(e^{Lh}p(h))
   =e^{Lh+aL^2}(e^{aD^2}p)(h+2aL)
\]
follows by expanding $(D+L)^2$, whose three terms commute. It also follows by completing the Gaussian square. With $V$ a standard real Gaussian it yields
\[
 \frac{f_{-\delta^2\tau}(b-2\delta^2\tau L+\delta w)}
 {\delta^2G(b)e^{\delta wL-\delta^2\tau L^2}}
 =J(\tau,w)
 =\mathbb E\bigl[(Y^2-1/4)R(\delta Y)\bigr],
 \quad Y=w+\sqrt{2\tau}V.
\]
All denominators are nonzero. The heat transform of an entire function of order at most one exists for every bounded complex heat parameter by Gaussian domination. The expectation above is even in $\sqrt\tau$ and therefore defines a holomorphic function of $\tau$ through zero.

Lemma~\ref{lem:centered-product} now gives, uniformly on fixed bounded complex $(\tau,w)$ sets,
\begin{align*}
 J(\tau,w)&=w^2+2\tau-\tfrac14-\tfrac\kappa2P_4(\tau,w)
             +\eta_3P_5(\tau,w)+O(\kappa^2),\\
 P_4(\tau,w)&=w^4+(12\tau-\tfrac14)w^2+12\tau^2-\tfrac12\tau,\\
 P_5(\tau,w)&=w^5+(20\tau-\tfrac14)w^3
                       +(60\tau^2-\tfrac32\tau)w,\\
 \eta_3&=\delta^3T_3/3,\qquad |\eta_3|\le\kappa^{3/2}/3.
\end{align*}
For completeness, the integrated remainder is bounded by a constant times
\[
 \kappa^2\mathbb E[(1+|V|)^6e^{C\kappa(1+V^2)}]=O(\kappa^2)
\]
for small absolute $\kappa$. Cauchy's formula gives the same remainder order for every fixed derivative on a smaller compact set. The correction of order $\kappa$ is even in $w$, and that of order $\kappa^{3/2}$ is odd.

For real $0\le\tau\le1/4$, the polynomial $w^2+2\tau-1/4$ has exactly two roots in $|w|<2$, and its modulus on the boundary is at least $15/4$. Rouch\'e's theorem therefore tracks exactly two roots of $J$ in the moving disk. Likewise $J_w=2w+O(\kappa)$ has exactly one zero there. Since $J_{ww}=2+O(\kappa)$, this critical point is a real analytic function $w_*(\tau)$. Evaluating $J_w$ at zero and using the parity just noted gives
$w_*(\tau)=O(\kappa^{3/2})$.

The critical value $D(\tau)=J(\tau,w_*(\tau))$ satisfies
\[
 D(\tau)=2\tau-\tfrac14
       -\tfrac\kappa2(12\tau^2-\tfrac12\tau)+O(\kappa^2),
 \qquad D'(\tau)=2+O(\kappa)>0.
\]
Its signs at $0$ and $1/4$ are opposite, so it has exactly one zero $\tau_c$. First $\tau_c=1/8+O(\kappa)$. At $\tau=1/8$ the expression $12\tau^2-\tau/2$ equals $1/8$, whence
$\tau_c=1/8+\kappa/32+O(\kappa^2)$.
At this point $J_{ww}\ne0$ and $D'\ne0$. The local quadratic factor thus has a simple discriminant zero. Positivity of $J_{ww}$ and $D'$ gives two real roots before the crossing and a conjugate pair after it. No further collision can occur because $J_w$ has only one zero and its critical value is strictly increasing. Returning to the original coordinates proves the result, including the spatial error. For $S=0$ all formulas reduce exactly to the heat transform of a quadratic times a linear exponential.
\end{proof}

The classical Lehmer-pair quantity of Csordas, Smith and Varga is
\[
 g=\sum_{\xi\ne b\pm\delta/2}
 \bigl((b-\delta/2-\xi)^{-2}+(b+\delta/2-\xi)^{-2}\bigr).
\]
Their Definition~1 and Theorem~1 \cite[pp.~109--110]{csv} use
$\delta^2g<4/5$ and give
$((1-5\delta^2g/4)^{4/5}-1)/(8g)\le\Lambda$.
Termwise symmetric expansion, justified by $\delta^2S\to0$, gives
$g=2S+O(\delta^2S^2)$. The first two terms of their bound are consequently
$-\delta^2/8-\delta^4S/32+O(\delta^6S^2)$, matching the negative of the collision time above. The theorem identifies the continued local pair and its unique crossing.

The two-angle example makes the normalization visible. For
$f(z)=2\cos(\delta/2)-2\cos(z-b)$ with $0<\delta<\pi$,
\[
 f_{-a}(z)=2\cos(\delta/2)-2e^{-a}\cos(z-b),\qquad
 a_c=-\log\cos(\delta/2)=\delta^2/8+\delta^4/192+O(\delta^6).
\]
The complementary periodic divisor has
$S=(2\sin^2(\delta/4))^{-1}-8/\delta^2=1/6+O(\delta^2)$, giving the same second coefficient.
